\begin{question}
\noindent A microbiologist is examining a circular petri dish of radius $6.5$ cm under a microscope. A straight scratch (a chord) on the base of the dish, $AB$, has length $10$ cm.

The point $O$ is the centre of the dish, and $M$ is the midpoint of $AB$.

\begin{center}
\begin{tikzpicture}[scale=0.9]
    \coordinate (O) at (0,0);
    \coordinate (A) at (-5,3.9);
    \coordinate (B) at (5,3.9);
    \coordinate (M) at (0,3.9);

    \draw[fill=lightYellow] (O) circle (6.5);
    \fill[black] (O) circle (2pt);
    \draw (A) -- (B);
    \draw (O) -- (A);
    \draw (O) -- (B);
    \draw (O) -- (M);

    \tkzMarkRightAngle[size=.25](O,M,B);

    \node[below] at (O) {$O$};
    \node[left] at (M) {$M$};
    \node[above left] at (A) {$A$};
    \node[above right] at (B) {$B$};

    \node[left] at (-2.5,2.1) {$6.5$ cm};
    \node[above] at (-2.5,3.9) {$5$ cm};
\end{tikzpicture}
\end{center}

\noindent Work out the distance $OM$, the perpendicular distance from the centre $O$ of the dish to the scratch $AB$.

\vspace{4.5cm}
\begin{flushright}
    \answerunit{cm}{3}
\end{flushright}
\end{question}
